%% ma: L10-B02-0006
%% lop: 10
%% chuong: 1
%% bai: 2
%% dang: 10.2.2
%% loai: TN4
%% muc: TH
%% dap_an: "D"
%% nguon: "Ảnh giáo viên gửi 10/10/2026 (ThS. Nguyễn Hoàng Việt, Chương 2 Tập hợp), A-Đề 1, Phần 1 Câu 4"
%% kien_thuc: "Bài 2 (tập hợp và các phép toán trên tập hợp)"
%% trang_thai: chinh-thuc
\begin{ex}
Cho tập hợp $A=\{x+1\mid x\in\mathbb N,\ x\le5\}$. Số phần tử của tập hợp $A$ là
\choice{$8$}{$7$}{$5$}{\True $6$}
\loigiai{$x\in\mathbb N$, $x\le5$ nên $x\in\{0;1;2;3;4;5\}$, suy ra $A=\{1;2;3;4;5;6\}$ có $6$ phần tử.}
\end{ex}
